\chapter{Einleitung}
\label{chap:einleitung}

\section{Motivation und Problemstellung}
\label{sec:motivation}

Mit europäischen Regelwerken wie der NIS-2-Richtlinie~\cite{NIS2Directive},
dem Cyber Resilience Act (CRA)~\cite{CyberResilienceAct} und der Digital
Operational Resilience Regulation (DORA)~\cite{DORARegulation} steigen die
Anforderungen an Organisationen, ihre IT-Systeme nicht nur präventiv abzusichern, sondern deren Widerstandsfähigkeit gegenüber Sicherheitsvorfällen ganzheitlich zu betrachten. Neben klassischen Schutzmaßnahmen rücken dabei insbesondere die Fähigkeiten zur Erkennung, Reaktion und Wiederherstellung nach Sicherheitsvorfällen sowie die Nachweisbarkeit dieser Prozesse in den Vordergrund. Resilienz wird damit zu einer Eigenschaft, die nicht nur konzipiert, sondern auch überprüft werden muss.

Für Forschung und Entwicklung ergibt sich daraus die Frage, wie sich solche Fähigkeiten
unter realistischen, zugleich aber kontrollierten Bedingungen untersuchen und bewerten
lassen. Produktive Systeme sind dafür nur eingeschränkt geeignet, da Sicherheitsvorfälle
dort erhebliche Risiken verursachen können. Gleichzeitig fehlt es an
standardisierten Testumgebungen, die reproduzierbare Experimente im Bereich der
Cyber-Resilienz erlauben \cite{Chouliaras2021}.

Cyber Ranges adressieren dieses Spannungsfeld, indem sie virtuelle, isolierte
IT-Landschaften bereitstellen, in denen Angriffs- und Verteidigungsszenarien realitätsnah
nachgebildet werden können. Während viele existierende Cyber-Range-Plattformen jedoch
primär für Schulungs- und Trainingszwecke entwickelt wurden, ist ihr systematischer
Einsatz als Forschungsumgebung -- insbesondere zur reproduzierbaren Evaluation von
Sicherheitswerkzeugen -- bislang nur teilweise etabliert \cite{IntroductionCyberRanges2022}.

Genau hier setzt die vorliegende Arbeit an. Um Sicherheits- und Analysewerkzeuge
objektiv miteinander vergleichen oder resilienzbezogene Eigenschaften eines Systems
beurteilen zu können, müssen identische Szenarien unter gleichen Randbedingungen
wiederholt ausführbar sein. Fehlt diese Reproduzierbarkeit, lassen sich beobachtete
Unterschiede nicht zweifelsfrei auf das untersuchte Werkzeug zurückführen. Die
Problemstellung besteht somit darin, eine geeignete Cyber-Range-Plattform so aufzubauen
und zu betreiben, dass sie als reproduzierbare Forschungsumgebung für
Cyber-Resilienz-Untersuchungen dienen kann, und dies anhand konkreter, wiederholbarer
Use Cases nachzuweisen.

\section{Zielsetzung}
\label{sec:zielsetzung}

Ziel dieser Arbeit ist es, eine geeignete Cyber-Range-Plattform auszuwählen, aufzubauen
und als reproduzierbare Forschungsumgebung für Fragestellungen der Cyber-Resilienz
nutzbar zu machen. Im Mittelpunkt stehen reproduzierbare Anwendungsszenarien (Use Cases),
an denen sich Sicherheits- und Analysewerkzeuge unter identischen Bedingungen einbinden
und vergleichen lassen sowie die Funktionsfähigkeit der Forschungsumgebung selbst belegen.
Der erste und im Detail ausgearbeitete Use Case ist ein Vergleich von Schwachstellen-Scannern
unterschiedlicher Charakteristik (klassisch, modern und kommerziell) gegen ein definiert
verwundbares Zielsystem.

Konkret verfolgt die Arbeit folgende Teilziele:
\begin{itemize}
  \item Es werden objektive Kriterien zur Bewertung von Cyber-Range-Plattformen aus
        Forschungsperspektive definiert.
  \item Auf dieser Basis wird die am besten geeignete Plattform ausgewählt und als
        Forschungsumgebung implementiert.
  \item Es werden reproduzierbare Cyber-Resilienz-Use-Cases konzipiert und technisch
        umgesetzt, in die externe Sicherheits- bzw.\ Analysewerkzeuge eingebunden werden.
  \item Es wird eine Evaluationsmethodik entwickelt, mit der sich Funktionsfähigkeit und
        Reproduzierbarkeit der Umgebung und der Use Cases überprüfen lassen.
\end{itemize}

Das angestrebte Ergebnis ist eine funktionsfähige, modular erweiterbare Forschungsumgebung,
die wiederholbare Experimente im Bereich der Cyber-Resilienz ermöglicht, gemeinsam mit einer
belastbaren Methodik zu deren Bewertung.

\section{Forschungsfragen}
\label{sec:forschungsfragen}

Aus der Problem- und Zielsetzung leitet sich die folgende übergeordnete Forschungsfrage ab:

\begin{quote}
Wie kann eine geeignete Cyber-Range-Plattform ausgewählt, aufgebaut und eingesetzt werden,
um eine funktionsfähige Forschungsumgebung für Cyber-Resilienz mit reproduzierbaren Use
Cases bereitzustellen?
\end{quote}

Zur Beantwortung der Forschungsfrage werden die folgenden Untersuchungsfragen betrachtet:

\begin{enumerate}
  \item Nach welchen objektiven Kriterien lassen sich Cyber-Range-Plattformen für Forschung
        und Tool-Integration bewerten?
  \item Welche Plattform erfüllt diese Anforderungen am besten und eignet sich für den Aufbau
        einer forschungsorientierten Umgebung?
  \item Wie können realistische und reproduzierbare Cyber-Resilienz-Use-Cases auf der
        ausgewählten Plattform modelliert und umgesetzt werden?
  \item Welche technischen Schnittstellen sind notwendig, um Sicherheits- oder
        Monitoring-Tools in die Cyber Range einzubinden?
  \item Wie lässt sich die Funktionsfähigkeit und Reproduzierbarkeit der entwickelten Use
        Cases evaluieren?
\end{enumerate}

\section{Methodik}
\label{sec:methodik}

Die Arbeit folgt einem \emph{Design-Science-Research}-Ansatz (DSR), der die Konstruktion
und Evaluation eines praktischen Artefakts -- hier der Forschungsumgebung samt Use Cases --
in den Mittelpunkt stellt und mit einer systematischen Literaturrecherche kombiniert
\cite{Hevner2004DesignScience}. Das Vorgehen orientiert sich an den etablierten Phasen
gestaltungsorientierter Forschung \cite{Peffers2007DSRM} und gliedert sich in:

\begin{enumerate}
  \item \textbf{Analyse:} Untersuchung existierender Cyber-Range-Plattformen sowie
        wissenschaftlicher Ansätze zu Szenarienmodellierung, Reproduzierbarkeit und
        Tool-Integration.
  \item \textbf{Plattformvergleich:} strukturierter Vergleich mehrerer Plattformen anhand
        definierter Kriterien und Auswahl der geeignetsten Lösung.
  \item \textbf{Konzeption:} Entwurf reproduzierbarer Cyber-Resilienz-Use-Cases inklusive
        Evaluationskriterien.
  \item \textbf{Implementierung:} technische Umsetzung der Use Cases auf der ausgewählten
        Plattform inklusive der notwendigen Integrationsschnittstellen.
  \item \textbf{Evaluation:} Überprüfung von Funktionsfähigkeit, Reproduzierbarkeit und
        Erweiterbarkeit der entwickelten Umgebung und Use Cases.
  \item \textbf{Reflexion:} Ableitung von Empfehlungen für den wissenschaftlichen Einsatz von
        Cyber Ranges.
\end{enumerate}

\section{Aufbau der Arbeit}
\label{sec:aufbau}

Kapitel~\ref{chap:einleitung} umreißt Motivation, Zielsetzung, Forschungsfragen und Methodik.
Kapitel~\ref{chap:grundlagen} erarbeitet die fachlichen Grundlagen zu Cyber Ranges, Cyber-Resilienz
und Schwachstellen-Scannern. Kapitel~\ref{chap:stand-der-technik} ordnet den Stand der Technik ein,
betrachtet bestehende Cyber-Range-Plattformen und leitet die Forschungslücke ab.
Kapitel~\ref{chap:poc} erprobt die beiden zentralen quelloffenen Kandidaten praktisch in einem Proof
of Concept. Aufbauend darauf entwickelt Kapitel~\ref{chap:auswahl} einen Kriterienkatalog, vergleicht
die Plattformen und begründet die Auswahl. Kapitel~\ref{chap:uc1} setzt den ersten Use Case auf der
gewählten Plattform um und führt von der Konzeption über die Evaluationsmethodik und die
Implementierung bis zu den Ergebnissen und deren Diskussion. Anhang~\ref{chap:quelltext-uc1} enthält
den vollständigen Quelltext der Use-Case-Umgebung.
